\chapter{矩阵}
\orig{6--23}

\section{矩阵乘法}

\textbf{定义} 设
\[
 A=(a_{ij})_{m\times n},\qquad B=(b_{ij})_{n\times p}.
\]
矩阵
\[
 C=(c_{ij})_{m\times p},
\]
其中
\[
 c_{ij}=a_{i1}b_{1j}+a_{i2}b_{2j}+\cdots+a_{in}b_{nj}
      =\sum_{k=1}^{n}a_{ik}b_{kj},
\]
称为 $A$ 与 $B$ 的乘积，记为
\[
 C=AB.
\]
注意，两个矩阵只有当第一个矩阵的列数等于第二个矩阵的行数时才能相乘。

矩阵乘法不适合交换律：
\begin{enumerate}[label=\roman*),leftmargin=2.4em,itemsep=0.35em]
\item $A_{m\times n}B_{n\times p}$ 有意义，当 $m\ne p$ 时，$B_{n\times p}A_{m\times n}$ 没有意义。例如 $A_{4\times3}B_{3\times5}$ 有意义，而 $B_{3\times5}A_{4\times3}$ 没有意义。
\item $A_{m\times n}B_{n\times m}$ 与 $B_{n\times m}A_{m\times n}$ 都有意义；当 $m\ne n$ 时，它们阶数不等。例如 $A_{4\times3}B_{3\times4}$ 是 4 阶矩阵，$B_{3\times4}A_{4\times3}$ 是 3 阶矩阵。
\item $A_{n\times n}B_{n\times n}$ 与 $B_{n\times n}A_{n\times n}$ 的阶数都是 $n$，也不一定相等。
\end{enumerate}

\textbf{例}
\[
 A=\begin{pmatrix}1&2\\2&1\end{pmatrix},\qquad
 B=\begin{pmatrix}2&-3\\3&1\end{pmatrix}.
\]
则
\[
 AB=\begin{pmatrix}8&-1\\7&-5\end{pmatrix},\qquad
 BA=\begin{pmatrix}-4&1\\5&7\end{pmatrix}.
\]

\section{有关矩阵乘法的例题}

\begin{enumerate}
\item \textbf{存在零因子。} 即 $A\ne O$，$B\ne O$，而有 $AB=O$。

例如
\[
 A=\begin{pmatrix}1&0\\0&0\end{pmatrix},\qquad
 B=\begin{pmatrix}0&0\\0&1\end{pmatrix},
\]
但
\[
 AB=\begin{pmatrix}0&0\\0&0\end{pmatrix}.
\]
一般地，对 $n\ge2$，可取
\[
 A=\operatorname{diag}(1,0,\ldots,0),\qquad
 B=\operatorname{diag}(0,1,0,\ldots,0),
\]
则 $A\ne O$，$B\ne O$，而 $AB=O$。

又如
\[
 A=\begin{pmatrix}1&-1\\-1&1\\1&-1\end{pmatrix},\qquad
 B=\begin{pmatrix}1&2\\1&2\end{pmatrix},
\]
而
\[
 AB=
 \begin{pmatrix}1&-1\\-1&1\\1&-1\end{pmatrix}
 \begin{pmatrix}1&2\\1&2\end{pmatrix}
 =\begin{pmatrix}0&0\\0&0\\0&0\end{pmatrix}.
\]

\item 上例中同阶矩阵相乘的情形里，$A\ne O$，$B\ne O$ 而有 $AB=O$，同时有 $BA=O$。现举一例是 $A\ne O$，$B\ne O$ 而有 $AB=O$，而 $BA\ne O$。

取
\[
 A=\begin{pmatrix}1&1\\0&0\end{pmatrix},\qquad
 B=\begin{pmatrix}1&1\\-1&-1\end{pmatrix}.
\]
那么
\[
 AB=\begin{pmatrix}0&0\\0&0\end{pmatrix},
\qquad
 BA=\begin{pmatrix}1&1\\-1&-1\end{pmatrix}.
\]
和上例一样，可以举出同样结论的 $n$ 阶矩阵的例。

\item \textbf{乘法消去律不成立。}
\[
 A\ne O,\qquad AB=AC
\]
未必有 $B=C$。

例如
\[
 A=\begin{pmatrix}1&1\\-1&-1\end{pmatrix},\quad
 B=\begin{pmatrix}1&-1\\-1&1\end{pmatrix},\quad
 C=\begin{pmatrix}-1&1\\1&-1\end{pmatrix}.
\]
显然 $A\ne O$，$AB=AC$，而 $B\ne C$。

\item 一般说来
\[
 (AB)^k\ne A^kB^k.
\]
例如
\[
 A=\begin{pmatrix}2&1\\3&2\end{pmatrix},\qquad
 B=\begin{pmatrix}1&-1\\1&1\end{pmatrix}.
\]
有
\[
 AB=\begin{pmatrix}3&-1\\5&-1\end{pmatrix},
\qquad
 (AB)^2=\begin{pmatrix}4&-2\\10&-4\end{pmatrix},
\]
而
\[
 A^2=\begin{pmatrix}7&4\\12&7\end{pmatrix},
\qquad
 B^2=\begin{pmatrix}0&-2\\2&0\end{pmatrix},
\]
从而
\[
 A^2B^2=\begin{pmatrix}8&-14\\14&-24\end{pmatrix},
\qquad
 (AB)^2\ne A^2B^2.
\]

$n$ 阶矩阵的例，只要把上述二阶矩阵置于左上角，并在右下角补上 $n-2$ 阶单位矩阵即可。明显地，只要 $AB=BA$，就有
\[
 (AB)^k=A^kB^k.
\]

\item 我们知道
\[
 (AB)'=B'A',
\]
可是未必有
\[
 (AB)'=A'B'.
\]
仍取
\[
 A=\begin{pmatrix}2&1\\3&2\end{pmatrix},\qquad
 B=\begin{pmatrix}1&-1\\1&1\end{pmatrix}.
\]
则
\[
 A'=\begin{pmatrix}2&3\\1&2\end{pmatrix},\qquad
 B'=\begin{pmatrix}1&1\\-1&1\end{pmatrix},
\]
\[
 AB=\begin{pmatrix}3&-1\\5&-1\end{pmatrix},\qquad
 A'B'=\begin{pmatrix}-1&5\\-1&3\end{pmatrix}.
\]
故有 $(AB)'\ne A'B'$。

\item
\[
 (AB)^{-1}=B^{-1}A^{-1},
\]
但未必有
\[
 (AB)^{-1}=A^{-1}B^{-1}.
\]
仍取
\[
 A=\begin{pmatrix}2&1\\3&2\end{pmatrix},\qquad
 B=\begin{pmatrix}1&-1\\1&1\end{pmatrix}.
\]
有
\[
 A^{-1}=\begin{pmatrix}2&-1\\-3&2\end{pmatrix},\qquad
 B^{-1}=\begin{pmatrix}\frac12&\frac12\\-\frac12&\frac12\end{pmatrix},
\]
\[
 AB=\begin{pmatrix}3&-1\\5&-1\end{pmatrix},\qquad
 (AB)^{-1}=\begin{pmatrix}-\frac12&\frac12\\-\frac52&\frac32\end{pmatrix},
\]
而
\[
 A^{-1}B^{-1}=\begin{pmatrix}\frac32&\frac12\\-\frac52&-\frac12\end{pmatrix}.
\]
故 $(AB)^{-1}\ne A^{-1}B^{-1}$。

\item 在本书默认的特征 $0$ 数域语境下，不存在 $n$ 阶矩阵 $A,B$ 满足
\[
 AB-BA=E.
\]
但可存在线性变换 $A,B$ 满足关系
\[
 AB-BA=E.
\]
请参看线性变换的有关例题。\jiaokan{若不限制底域特征，“不存在 $n$ 阶矩阵 $A,B$ 满足 $AB-BA=E$”不能无条件成立。在特征 $0$（以及更一般地，域的特征不整除 $n$）时，由 $\operatorname{tr}(AB-BA)=0$ 与 $\operatorname{tr}E=n$ 即得矛盾；此处据此补明语境。}

\item 一般说来，对 $n$ 阶矩阵 $A$ 和 $B$，等式
\[
 |A+B|=|A|+|B|
\]
不成立。

例如取
\[
 A=\operatorname{diag}(1,\ldots,1,0),\qquad
 B=\operatorname{diag}(0,\ldots,0,1).
\]
显然 $A+B=E$，因而
\[
 |A+B|=|E|=1,
\]
又
\[
 |A|+|B|=0+0=0.
\]

\item $n$ 阶矩阵 $A,B$，且 $A^2=E$、$B^2=E$，未必有 $(AB)^2=E$。

例如当 $n=2$ 时，取
\[
 A=\begin{pmatrix}1&2\\0&-1\end{pmatrix},\qquad
 B=\begin{pmatrix}1&3\\0&-1\end{pmatrix}.
\]
有
\[
 A^2=E,\qquad B^2=E,
\]
\[
 AB=\begin{pmatrix}1&1\\0&1\end{pmatrix},
\]
而
\[
 (AB)^m=\begin{pmatrix}1&m\\0&1\end{pmatrix}\qquad(m\text{ 为正整数}),
\]
所以 $(AB)^m\ne E$，特别地 $(AB)^2\ne E$。

\item 已知有理数域 $\Q$ 上 3 阶矩阵
\[
 A=\begin{pmatrix}
 2&0&0\\
 -3&1&\frac32\\
 2&\frac23&1
 \end{pmatrix}
\]
的行列式是 0，找一个 $\Q$ 上的非零 3 阶矩阵 $B$ 使
\[
 AB=BA=O.
\]
这样的 $B$ 是不是唯一的？

根据条件 $AB=BA=O$，找到
\[
 B=\begin{pmatrix}
 0&0&0\\
 -\frac92k&-\frac32k&\frac94k\\
 3k&k&-\frac32k
 \end{pmatrix},
\qquad k\ne0,\quad k\in\Q.
\]
确实 $AB=BA=O$。取 $k$ 的不同值，就得到不同的 $B$。
\end{enumerate}

\section{（实）对称阵}

\begin{enumerate}
\item 对称阵之和仍为对称阵；但对称阵之积未必是对称阵。

例如
\[
 A=\begin{pmatrix}1&3\\3&2\end{pmatrix},\qquad
 B=\begin{pmatrix}2&1\\1&1\end{pmatrix},
\]
而
\[
 AB=\begin{pmatrix}5&4\\8&5\end{pmatrix}
\]
不是对称阵。应当注意的是，1 阶对称阵之积仍为对称阵。

\item 设 $A$ 是（实）对称阵，如果 $A^2=O$，就有 $A=O$。我们说，$A$ 是对称阵的条件不容忽视。

例如
\[
 A=\begin{pmatrix}a&a\\-a&-a\end{pmatrix},\qquad a\in\R.
\]
显然 $A^2=O$；当 $a\ne0$ 时，$A\ne O$，此时 $A$ 不是对称阵。

$n$ 阶矩阵的情形，可令
\[
 A=\begin{pmatrix}
 a&a&0&\cdots&0\\
 -a&-a&0&\cdots&0\\
 0&0&0&&0\\
 \vdots&\vdots&&\ddots&\\
 0&0&0&&0
 \end{pmatrix},
\]
同样有 $A^2=O$ 而 $A\ne O$（$a\ne0$）。

\item 实对称阵的特征值都是实数，但特征值都是实数的实矩阵未必对称。

例如
\[
 A=\begin{pmatrix}4&4\\1&4\end{pmatrix}
\]
的特征值为 2 和 6，但 $A$ 不对称。

\item 实对称阵和对角阵相似，但是和对角阵相似的矩阵未必对称。

例如
\[
 A=\begin{pmatrix}1&0\\0&2\end{pmatrix},\qquad
 B=\begin{pmatrix}1&-\frac12\\0&2\end{pmatrix},
\]
取
\[
 T=\begin{pmatrix}2&1\\0&3\end{pmatrix},\qquad
 T^{-1}=\begin{pmatrix}\frac12&-\frac16\\0&\frac13\end{pmatrix}.
\]
有
\[
 B=T^{-1}AT.
\]
即 $A$ 与 $B$ 相似，$A$ 是对角阵，而 $B$ 不是对称阵。

\item $n$ 阶实对称阵有 $n$ 个线性无关的特征向量。反之，有 $n$ 个线性无关的特征向量不一定是实对称阵。

例如
\[
 A=\begin{pmatrix}
 4&6&0\\
 -3&-5&0\\
 -3&-6&1
 \end{pmatrix}.
\]
$A$ 的特征值为 $-2$ 和 $1$。属于 $-2$ 的特征向量是 $(-1,1,1)$；属于 $1$ 的特征向量是 $(-2,1,0)$、$(0,0,1)$。这三个特征向量线性无关，但 $A$ 不是对称阵。
\end{enumerate}

\section{反对称阵}

所谓反对称矩阵，即要求
\[
 A=-A'.
\]

\begin{enumerate}
\item 关于反对称阵之积。当阶数为 2 时，
\[
 A=\begin{pmatrix}0&a\\-a&0\end{pmatrix},\qquad
 B=\begin{pmatrix}0&-b\\b&0\end{pmatrix}
\]
均为反对称阵，而
\[
 AB=\begin{pmatrix}ab&0\\0&ab\end{pmatrix}.
\]
当 $a\ne0$ 且 $b\ne0$ 时，$AB$ 是对称阵，且为对角形阵，亦为纯量矩阵，不是反对称阵。

我们不能将这一结论推广到 $n$ 阶阵。例如
\[
 A=\begin{pmatrix}
 0&-1&-2\\1&0&-3\\2&3&0
 \end{pmatrix},\qquad
 B=\begin{pmatrix}
 0&-4&-5\\4&0&-6\\5&6&0
 \end{pmatrix},
\]
则
\[
 AB=\begin{pmatrix}
 -14&-12&6\\
 -15&-22&-5\\
 12&-8&-28
 \end{pmatrix},
\]
既不是对称阵也不是反对称阵。

又例如
\[
 A=\begin{pmatrix}
 0&0&2&0\\
 0&0&3&0\\
 -2&-3&0&0\\
 0&0&0&0
 \end{pmatrix},\qquad
 B=\begin{pmatrix}
 0&0&-3&0\\
 0&0&-1&0\\
 3&1&0&0\\
 0&0&0&0
 \end{pmatrix},
\]
都是反对称阵，而
\[
 AB=\begin{pmatrix}
 6&2&0&0\\
 9&3&0&0\\
 0&0&9&0\\
 0&0&0&0
 \end{pmatrix}
\]
既不是对称阵也不是反对称阵。

当然可以把以上两个例题推广到一般 $n$ 阶阵上去。同样，不宜忽略 1 阶的情形：
\[
 A=(0),\qquad B=(0),\qquad AB=(0),
\]
既可谓对称阵亦可谓反对称阵。

\item $A$ 是 $n$ 阶阵：当 $A$ 是对称阵时，$A^m$ 仍为对称阵；当 $A$ 是反对称阵时，$A^m$ 在 $m$ 为奇数时是反对称阵，在 $m$ 为偶数时是对称阵。

\item 不存在奇数阶的可逆反对称阵，也就是说，奇数阶的反对称阵的行列式为 0；而偶数阶的反对称阵的行列式不一定为 0。

例如
\[
 A=\begin{pmatrix}
 0&-1&0&0\\
 1&0&0&0\\
 0&0&0&-1\\
 0&0&1&0
 \end{pmatrix}
\]
显然 $|A|\ne0$。

当然偶数阶的反对称阵的行列式亦可为 0。例如
\[
 A=\begin{pmatrix}
 0&-1&0&0\\
 1&0&0&0\\
 0&0&0&0\\
 0&0&0&0
 \end{pmatrix},
\]
显然 $|A|=0$。

但是也有一个必然结论的情况，那就是 2 阶非零的反对称阵的行列式一定不为 0。例如
\[
 A=\begin{pmatrix}0&a\\-a&0\end{pmatrix},\qquad a\ne0,
\]
而
\[
 |A|=a^2\ne0.
\]
\end{enumerate}

\section{正定阵}

\begin{enumerate}
\item 正定阵的和还是正定阵，但正定阵的差未必是正定阵。

例如
\[
 A=\begin{pmatrix}1&1\\1&2\end{pmatrix},\qquad
 B=\begin{pmatrix}1&-1\\-1&2\end{pmatrix}
\]
都是正定阵，但
\[
 A-B=\begin{pmatrix}0&2\\2&0\end{pmatrix}
\]
不是正定阵。

\item 正定阵的积未必是正定阵。

\textbf{例 1}
\[
 A=\begin{pmatrix}1&1\\1&2\end{pmatrix},\qquad
 B=\begin{pmatrix}1&-1\\-1&2\end{pmatrix},
\]
而
\[
 AB=\begin{pmatrix}0&1\\-1&3\end{pmatrix}
\]
不是正定阵。

\textbf{例 2}
\[
 A=\begin{pmatrix}1&2\\2&5\end{pmatrix},\qquad
 B=\begin{pmatrix}1&-1\\-1&2\end{pmatrix}
\]
都是正定阵，而
\[
 AB=\begin{pmatrix}-1&3\\-3&8\end{pmatrix}
\]
不是正定阵。

还应指出的是，在 1 阶阵中，若 $A=(a)$、$B=(b)$ 都是正定阵，则其乘积
\[
 AB=(ab)
\]
也是正定阵。

\item $A$ 是正定阵，则 $A$ 的行列式 $|A|>0$。但对称阵 $B$ 的行列式 $|B|>0$，$B$ 未必是正定阵。

例如
\[
 B=\begin{pmatrix}1&0&0\\0&-1&0\\0&0&-1\end{pmatrix},
\]
而 $|B|=1>0$，但 $B$ 不是正定阵。

\item $A$ 是正定阵，则 $A$ 的主对角线上元素都大于零。但反之不真。

例如
\[
 A=\begin{pmatrix}1&3\\3&2\end{pmatrix},\qquad
 B=\begin{pmatrix}1&2\\2&4\end{pmatrix}
\]
的主对角元都大于零，但二者都不是正定阵。
\end{enumerate}

\section{正交阵}

所谓正交阵指的是 $n$ 阶实数矩阵 $A$ 满足条件
\[
 A'A=E.
\]
我们知道，正交阵之积仍为正交阵，那么正交阵之和是不是正交阵？

取以下两个 $n$ 阶阵
\[
 A=\operatorname{diag}(1,-1,1,\ldots,1),\qquad
 B=\operatorname{diag}(1,-1,-1,1,\ldots,1),
\]
都是正交阵，因为 $A'A=E$、$B'B=E$。但 $A+B$ 的对角元中出现 $2,-2,0,2,\ldots,2$，故
\[
 (A+B)'(A+B)\ne E.
\]
所以，正交阵的和不是正交阵。

另一个问题是：若 $A$ 是正交阵，则 $|A|=\pm1$，但反之不真。

例如
\[
 A=\begin{pmatrix}1&1\\1&2\end{pmatrix},
\qquad |A|=1,
\]
但
\[
 A'A=\begin{pmatrix}2&3\\3&5\end{pmatrix}\ne E.
\]
又取
\[
 B=\begin{pmatrix}1&1\\-1&-2\end{pmatrix},
\qquad |B|=-1,
\]
而
\[
 B'B=\begin{pmatrix}2&3\\3&5\end{pmatrix}\ne E.
\]
即 $A,B$ 都不是正交阵。

\section{等价矩阵、合同矩阵、相似矩阵}

\begin{enumerate}
\item 合同矩阵一定是等价矩阵，但反之不真。

取
\[
 B=\begin{pmatrix}1&1\\0&1\end{pmatrix},\qquad
 A=\begin{pmatrix}1&0\\0&1\end{pmatrix},
\]
\[
 P=\begin{pmatrix}1&0\\0&1\end{pmatrix},\qquad
 Q=\begin{pmatrix}1&1\\0&1\end{pmatrix}.
\]
显然有
\[
 B=PAQ,
\]
即 $A$ 与 $B$ 等价。但 $A$ 与 $B$ 不合同，否则
\[
 B=C'AC=C'C.
\]
令
\[
 C'=\begin{pmatrix}a&b\\c&d\end{pmatrix},\qquad
 C=\begin{pmatrix}a&c\\b&d\end{pmatrix},
\]
则
\[
 C'C=\begin{pmatrix}
 a^2+b^2&ac+bd\\
 ac+bd&c^2+d^2
 \end{pmatrix}=B.
\]
于是同一个量 $ac+bd$ 必须同时等于 $1$ 与 $0$，矛盾。

\item 相似矩阵一定是等价矩阵，但反之不真。

仍取
\[
 B=\begin{pmatrix}1&1\\0&1\end{pmatrix},\qquad
 A=E=\begin{pmatrix}1&0\\0&1\end{pmatrix}.
\]
知 $A$ 与 $B$ 等价，但 $A$ 与 $B$ 不相似，否则
\[
 B=T^{-1}AT=T^{-1}T=E,
\]
矛盾。

\item 相似矩阵未必合同。

例如
\[
 A=\begin{pmatrix}1&0\\0&2\end{pmatrix},\qquad
 B=\begin{pmatrix}1&-\frac12\\0&2\end{pmatrix},
\]
\[
 T=\begin{pmatrix}2&1\\0&3\end{pmatrix},\qquad
 T^{-1}=\begin{pmatrix}\frac12&-\frac16\\0&\frac13\end{pmatrix}.
\]
那么
\[
 B=T^{-1}AT,
\]
即 $A$ 与 $B$ 相似。但 $A$ 与 $B$ 不合同，否则
\[
 B=C'AC.
\]
令
\[
 C'=\begin{pmatrix}a&b\\c&d\end{pmatrix},\qquad
 C=\begin{pmatrix}a&c\\b&d\end{pmatrix},
\]
则
\[
 C'AC=\begin{pmatrix}
 a^2+2b^2&ac+2bd\\
 ac+2bd&c^2+2d^2
 \end{pmatrix}.
\]
比较非对角元，须有
\[
 ac+2bd=-\frac12,
\qquad
 ac+2bd=0,
\]
矛盾。

\item 合同矩阵未必相似。

取
\[
 B=\begin{pmatrix}1&0\\0&4\end{pmatrix},\qquad
 A=\begin{pmatrix}1&0\\0&1\end{pmatrix},
\]
\[
 C=\begin{pmatrix}1&0\\0&2\end{pmatrix},\qquad
 C'=\begin{pmatrix}1&0\\0&2\end{pmatrix}.
\]
则
\[
 B=C'AC.
\]
即 $A$ 与 $B$ 合同，但 $A$ 与 $B$ 不相似，否则
\[
 B=T^{-1}AT=T^{-1}T=E,
\]
矛盾。

\item $A$ 可逆，则有 $AB$ 与 $BA$ 相似，反之不真。

例如
\[
 A=\begin{pmatrix}1&0\\0&0\end{pmatrix},\qquad
 B=\begin{pmatrix}2&0\\0&0\end{pmatrix}.
\]
显然有 $AB=BA$，故 $AB$ 与 $BA$ 相似，而 $A$ 没有逆。

\item 相似矩阵有相同的特征多项式，但反之不真。

例如
\[
 A=\begin{pmatrix}1&0\\0&1\end{pmatrix},\qquad
 B=\begin{pmatrix}1&1\\0&1\end{pmatrix}.
\]
有
\[
 |\lambda E-A|=(\lambda-1)^2,
\qquad
 |\lambda E-B|=(\lambda-1)^2,
\]
即有相同特征多项式，可是 $A$ 与 $B$ 不相似。

\item 若
\[
 B=T^{-1}AT,
\qquad
 B=Q^{-1}AQ,
\]
则 $T$ 未必就等于 $Q$。

例如
\[
 A=B=\begin{pmatrix}1&1\\0&1\end{pmatrix},
\]
取
\[
 T=\begin{pmatrix}1&0\\0&1\end{pmatrix},\qquad
 Q=\begin{pmatrix}2&0\\0&2\end{pmatrix}.
\]
则有
\[
 B=T^{-1}AT,
\qquad
 B=Q^{-1}AQ,
\]
而 $T\ne Q$。
\end{enumerate}
